% ===== SECTION 5: COST CONTROLS =====
\section{\texorpdfstring{\color{MidnightBlue}Cost Controls}{Cost Controls}}

\begin{sidebar}{RFI Reference}
Section C: Questions to Field --- Cost Controls (Questions 1--4)\\[0.3cm]
\textit{Note: Per the RFI instructions, no pricing information is included. This section describes methodologies only.}
\end{sidebar}

% ----- Question 1 -----
\subsection{Question 1: Cost Per Bus Per Day}
\textit{``How do you address and apply cost per bus per day? (no pricing info, but explain your methodology)''}

\vspace{0.3cm}

Sentry's analytics platform calculates cost per bus per day as a comprehensive management metric---not a billing rate, but a tool for identifying inefficiency, benchmarking across zones, and informing procurement decisions. The methodology captures all variable and fixed costs attributable to each vehicle on each operating day, enabling RIDE to compare the true cost of service delivery across operators, vehicle types, and zones on an apples-to-apples basis.

\begin{description}[style=nextline, leftmargin=0.5cm]
    \item[\faIcon{calculator} \textbf{Cost Components:}] The per-bus-per-day metric aggregates seven cost components: (1)~driver wages and benefits (including RI prevailing wage obligations, Teamsters pension/healthcare contributions, and paid holiday costs); (2)~fuel consumption (GPS-verified actual gallons, not estimated); (3)~scheduled maintenance (prorated PM costs by vehicle age and interval); (4)~unscheduled maintenance (road calls, breakdowns, warranty-excluded repairs); (5)~insurance premiums allocated per vehicle-day; (6)~vehicle depreciation or lease cost amortized over operating days; and (7)~administrative overhead (supervision, dispatch, facility costs allocated proportionally). Each component is tracked separately so RIDE can identify which cost element is driving changes---for example, distinguishing fuel price increases (external) from rising maintenance costs (fleet age issue).

    \item[\faIcon{layer-group} \textbf{Vehicle Type Stratification:}] Our methodology stratifies cost per bus per day by vehicle type, recognizing that the cost profile of a 10-passenger van differs significantly from a Type A wheelchair-accessible bus or a Type C/D full-size bus. With 147 vans, 93 Type~A buses, 30 Type~A wheelchair buses, and 64 Type~B/C/D buses in RIDE's fleet, a single blended cost-per-bus-per-day figure would obscure meaningful variation. Stratification reveals, for example, that a 10-passenger van costs substantially less per day than a Type~C 72-passenger bus---and when a route carries only 3--5 students (common in the RIDE system, where Zone~D averages one student per route), the platform flags the opportunity to right-size the vehicle from a bus to a van, capturing the \textasciitilde\$26,000 per-vehicle savings the Study Commission identified. This analysis directly supports RIDE's stated goal of maximizing van usage under the new 10-passenger vehicle law.

    \item[\faIcon{chart-line} \textbf{Trend Analysis:}] The platform maintains a rolling 12-month trend analysis of cost per bus per day, segmented by zone, operator, and vehicle type. This longitudinal view identifies cost creep before it compounds---a 2\% monthly increase in Zone~C maintenance costs, for instance, signals fleet aging that warrants vehicle replacement planning. Seasonal variation is baselined (summer school sessions have different cost profiles than the academic year), and the impact of specific efficiency initiatives is isolated: when a van conversion reduces the cost per bus per day on a given route by 15\%, that savings is tracked and verified against the projection. RIDE can use these trends to set performance-based contract incentives tied to measurable cost reduction.
\end{description}

% ----- Question 2 -----
\subsection{Question 2: Cost Per Mile}
\textit{``What is your company's cost per mile methodology?''}

\vspace{0.3cm}

Sentry's cost-per-mile methodology uses GPS-verified mileage data---not odometer readings or estimated distances---to calculate precise, auditable cost metrics. By capturing actual miles driven via integrated GPS tracking hardware, the platform eliminates the discrepancies between estimated route miles and real-world driving patterns that can distort cost analysis by 10--15\%.

\begin{description}[style=nextline, leftmargin=0.5cm]
    \item[\faIcon{road} \textbf{Total Cost Per Mile:}] Total cost per mile divides all-in operating costs (driver labor, fuel, maintenance, insurance, depreciation, overhead) by total revenue miles---miles driven with students aboard. By excluding deadhead miles from the denominator, this metric reveals the true cost of productive service. In RIDE's system, where Zones~C and~D have deadhead ratios of 36--51\%, the total cost per \textit{revenue} mile is dramatically higher than the cost per \textit{total} mile, exposing the real economic impact of routing inefficiency.

    \item[\faIcon{gas-pump} \textbf{Fuel Cost Per Mile:}] Fuel cost per mile is tracked as a standalone metric to isolate fuel price volatility (an external factor RIDE and operators cannot control) from operational efficiency (which they can). GPS-verified mileage eliminates odometer discrepancies and enables per-vehicle fuel economy monitoring. The platform flags vehicles with fuel consumption more than 15\% above their class average---indicating mechanical issues (injector problems, tire pressure, drag) that increase costs and emissions simultaneously.

    \item[\faIcon{wrench} \textbf{Maintenance Cost Per Mile:}] Maintenance cost per mile is tracked by vehicle age, type, and maintenance category (scheduled PM vs. unscheduled breakdown). The platform calculates a rolling 6-month maintenance cost per mile for each vehicle and compares it against replacement cost thresholds. When a vehicle's maintenance cost per mile exceeds 50\% of its replacement lease cost per mile, the system flags it as approaching end-of-economic-life---enabling RIDE to plan fleet replacement proactively rather than reactively after a breakdown strands students.

    \item[\faIcon{divide} \textbf{Loaded vs. Deadhead Mile Distinction:}] This is the most critical distinction in cost-per-mile analysis. Cost per \textit{loaded} mile (with students aboard) is the true efficiency metric; cost per \textit{total} mile (including deadhead) masks inefficiency. Sentry's platform separates these automatically via GPS geofencing at school locations and student pickup/drop-off points. When a vehicle crosses a geofence boundary at a school or designated stop, the system timestamps the transition between loaded and deadhead status. For RIDE, this means every mile is categorized and costed accurately---and the ``expenditure follows the child'' billing methodology can be validated against actual loaded miles rather than estimated route distances.
\end{description}

\begin{sidebar}{Cost Transparency for RIDE}
RIDE currently uses an ``expenditure follows the child'' billing methodology, where the total cost of a route is proportionally divided among LEAs based on the number of students each district has on that route. Today, this billing process generates approximately 700 checks per year from districts to RIDE---an outdated manual system that RIDE's own October~2024 overview identifies as a weakness. Sentry's platform automates this cost allocation in real-time, tying GPS-verified route data directly to per-student billing and eliminating manual reconciliation. As an independent technology provider, Sentry gives RIDE visibility into cost metrics across all operators and zones---eliminating the information asymmetry that exists when operators self-report their own cost data.
\end{sidebar}

% ----- Question 3 -----
\subsection{Question 3: Spare Ratio Philosophy}
\textit{``What is your company's spare ratio philosophy?''}

\vspace{0.3cm}

Spare vehicle management is an area where technology-driven analytics can yield significant cost savings. Each spare vehicle represents capital tied up in a non-revenue asset---purchase or lease cost, insurance, registration, and parking space. Sentry's approach uses data to right-size the spare pool: maintaining enough spares to ensure uninterrupted service while avoiding the excess that inflates RIDE's per-route costs.

\begin{description}[style=nextline, leftmargin=0.5cm]
    \item[\faIcon{percent} \textbf{Target Spare Ratio:}] The industry standard spare ratio for school bus fleets is 15--20\%. Sentry recommends a tiered target: 12--15\% for fleets with average vehicle age under 5 years (where warranty coverage reduces unscheduled downtime), and 18--20\% for older fleets or mixed fleets with high wheelchair-accessible vehicle counts (which have more complex mechanical systems). For RIDE's 334-vehicle fleet, this translates to approximately 50--67 spare vehicles system-wide. The ratio should be calculated per zone to account for local fleet age and condition---Zone~D's older fleet may require a higher local spare ratio than Zone~A1.

    \item[\faIcon{sync-alt} \textbf{Dynamic Spare Deployment:}] The platform tracks daily spare deployment by vehicle ID, recording which spare replaced which primary vehicle, the reason for substitution (PM, breakdown, driver assignment, accident), and duration of deployment. Over time, pattern analysis reveals actionable insights: a specific route that pulls spares three times per month may be too demanding for its assigned vehicle type; a vehicle class with rising spare utilization signals systematic reliability decline. When spare utilization trends upward across a zone for three or more consecutive weeks, the system generates an alert to RIDE and the zone operator---this is a leading indicator of fleet health problems that, if unaddressed, will escalate to service disruptions.

    \item[\faIcon{balance-scale} \textbf{Right-Sizing the Spare Pool:}] Sentry's analytics engine models the relationship between PM compliance, mean distance between failures, and required spare count. As PM compliance improves (target $\geq$ 98\%) and MDBF increases, the probability of unscheduled vehicle downtime decreases---and fewer spares are needed to maintain the same service reliability level. The platform runs monthly simulations: if Zone~B's MDBF improves from 10,000 to 14,000 miles over six months, the model may recommend reducing that zone's spare pool from 7 vehicles to 5---freeing two vehicles for reassignment or retirement and saving RIDE the associated insurance, parking, and depreciation costs. This data-driven approach replaces the common practice of maintaining a fixed spare ratio regardless of actual fleet performance.
\end{description}

% ----- Question 4 -----
\subsection{Question 4: Average Driver Hours Per Route}
\textit{``How do you address and apply average driver hours per route?''}

\vspace{0.3cm}

Average driver hours per route is the primary labor cost metric for student transportation, and RIDE's published data shows dramatic variation across zones---from 4h28m (Zone~A3) to 7h05m (Zone~D) including deadhead. Sentry's platform decomposes this aggregate into its component parts, enabling targeted optimization of each segment rather than treating the total as a monolithic number.

\begin{description}[style=nextline, leftmargin=0.5cm]
    \item[\faIcon{clock} \textbf{Components of Driver Hours:}] The platform captures six distinct components of driver hours per route: (1)~pre-trip inspection time (from driver app login to vehicle departure), (2)~deadhead to first pickup (GPS-tracked empty miles from depot to first student), (3)~loaded service time (first student pickup to last student drop-off), (4)~deadhead return (last drop-off to depot), (5)~post-trip inspection (vehicle check and interior sweep), and (6)~administrative time (fueling, paperwork, training). Each component is timestamped automatically---the driver app records inspection start/end, GPS geofences capture pickup and drop-off transitions, and the dispatch system logs administrative activities. This decomposition reveals, for example, that Zone~D's 7h05m average includes 3h38m of pure deadhead---meaning only 49\% of paid driver time is spent transporting students.

    \item[\faIcon{compress-arrows-alt} \textbf{Optimization Levers:}] Each component has a distinct optimization lever: deadhead time is reduced through the routing optimization, dynamic depot assignment, and cross-zone strategies described in Section~4. Pre-trip and post-trip inspection time is reduced through digital checklists on tablets---eliminating paper forms and enabling parallel data entry during the physical walk-around (our operators average 8--10 minutes per digital inspection vs.\ 15--20 minutes on paper). Administrative time is reduced through automated timekeeping (GPS-based clock-in/clock-out eliminates manual time cards), electronic fueling records, and digital route sheets that replace printed manifests. Combined, these optimizations can reduce non-service driver time by 20--30 minutes per route per day---modest per route, but across 334 daily routes, that translates to approximately 110--165 fewer paid driver-hours per day system-wide.

    \item[\faIcon{users} \textbf{Split-Shift vs. Straight-Time Analysis:}] The platform models both split-shift and straight-time scenarios for every route and zone, calculating the total labor cost under each approach. In a split-shift model, the driver works the AM run, goes off-clock for 3--4 hours, and returns for the PM run---reducing paid hours but potentially increasing driver dissatisfaction and turnover. In a straight-time model, the driver remains on-clock and covers midday routes (PCCT transfers, McKinney-Vento trips, field trips)---increasing paid hours but eliminating a second deadhead cycle and improving vehicle utilization. For RIDE, the optimal answer varies by zone: Zone~A3 (average 4h28m total, high student density) may favor split shifts, while Zone~D (average 7h05m, remote geography) may be more cost-effective as straight-time with midday route chaining. The platform runs both models and presents RIDE with a zone-by-zone cost comparison to inform contract structuring.
\end{description}

\mybox{
\textbf{\faIcon{info-circle} Note on RIDE's Current Data:} The RFI data shows average route times ranging from 4h28m (Zone~A3) to 7h05m (Zone~D) including deadhead. Sentry's platform would decompose these averages into their component parts, enabling RIDE to identify exactly where time (and money) is being lost on each route. Critically, this analysis must account for the system's changing demand profile: PCCT students---who constitute 50.87\% of riders---are subject to regional boundaries, while Special Education (24.29\%) and Homeless/Foster Care (24.84\%) riders are federally mandated regardless of region. A driver hour optimization that works for a predictable PCCT route may not apply to a McKinney-Vento route where student addresses change frequently.
}{MidnightBlue}
